% ============================================================================
% preamble.tex — форматиране по изискванията на ФКСТ.
%
% Всяка секция тук съответства на правило от REQUIREMENTS.md (Част А) и на
% реда му в таблицата на превода (Част Б). Не пипай стойностите без да
% провериш срещу REQUIREMENTS.md — те са нормативни, не са вкус.
%
% Двигател: pdfLaTeX (T2A + tempora). НЕ XeLaTeX — вж. REQUIREMENTS.md Б.0.
% ============================================================================

\usepackage{cmap}
\usepackage[T2A]{fontenc}
\usepackage[utf8]{inputenc}
\usepackage[bulgarian]{babel}

% --- Типография (А.4) ---------------------------------------------------
% tempora = Times-метрично съвместим шрифт с пълна кирилица под T2A.
\usepackage{tempora}
\usepackage[top=2.5cm, bottom=2.5cm, left=3.0cm, right=2.0cm, a4paper]{geometry}
\usepackage{setspace}
\setstretch{1.5}                     % „междуредие 1.5 lines“ — калибруем, вж. Б.2
% „без пренасяния на срички“ (А.4) — чрез ЗАБРАНИТЕЛНИ НАКАЗАНИЯ, не чрез
% \usepackage[none]{hyphenat}. Разликата е съществена: `hyphenat[none]` задава
% \hyphenchar=-1 и така изключва самия механизъм за прекъсване, от който зависи
% `breaklines` на listings — дългите редове в листингите тогава НЕ се пренасят и
% излизат извън полето (1227 препълвания дойдоха оттам). Наказанията постигат същия
% резултат за текста (TeX никога не пренася), но оставят механизма достъпен.
\usepackage[expansion=false]{microtype}
\DisableLigatures{encoding=*,family=tt*}
\hyphenpenalty=10000
\exhyphenpenalty=10000
\emergencystretch=1.5em              % цената на justified-без-преноси, вж. Б.4
\sloppy
\usepackage{indentfirst}             % отстъп и на ПЪРВИЯ абзац след заглавие, вж. Б.3
\setlength{\parindent}{1.25cm}
\setlength{\parskip}{0pt}

% --- Номерация на страниците: долу вдясно (А.4) -------------------------
\usepackage{fancyhdr}
\pagestyle{fancy}
\fancyhf{}
\renewcommand{\headrulewidth}{0pt}
\fancyfoot[R]{\thepage}

% --- Бележки под линия: 10 pt, single, justified (А.4) -------------------
\usepackage[hang,flushmargin]{footmisc}
\usepackage{ragged2e}
\renewcommand{\footnotelayout}{\fontsize{10}{12}\selectfont\setstretch{1.0}\justifying}

% --- Заглавия на раздели / параграфи / подпараграфи (А.7) ---------------
\usepackage{titlesec}
\newcommand{\sectionbreak}{\clearpage}   % „всеки раздел започва на нова страница“

\renewcommand{\thesection}{\Roman{section}}
\renewcommand{\thesubsection}{\arabic{section}.\arabic{subsection}}
\renewcommand{\thesubsubsection}{\arabic{section}.\arabic{subsection}.\arabic{subsubsection}}

% Раздел: римски, ALL CAPS, центриран, bold 14 pt
\titleformat{\section}[block]
  {\fontsize{14}{16.8}\bfseries\centering}{\thesection. }{0pt}{\MakeUppercase}

% Параграф: арабски, bold 14 pt, отстъп 1.25 cm
\titleformat{\subsection}[block]
  {\fontsize{14}{16.8}\bfseries}{\thesubsection. }{0pt}{}
\titlespacing*{\subsection}{1.25cm}{*3}{*2}

% Подпараграф: bold 13 pt, отстъп 2.0 cm
\titleformat{\subsubsection}[block]
  {\fontsize{13}{15.6}\bfseries}{\thesubsubsection. }{0pt}{}
\titlespacing*{\subsubsection}{2.0cm}{*2}{*1}

% --- Съдържание: разделите с ГЛАВНИ букви, 12 pt, bold (А.6) -------------
\usepackage{etoolbox}
\usepackage[titles]{tocloft}
% NB: кръпката ТРЯБВА да е СЛЕД tocloft — пакетът предефинира \l@section и
% иначе я изтрива (тихо: няма грешка, просто разделите излизат с малки букви).
\makeatletter
\patchcmd{\l@section}{#1}{\MakeUppercase{#1}}{}{}
\makeatother
\renewcommand{\cftsecleader}{\cftdotfill{\cftdotsep}}
\renewcommand{\cftsubsecleader}{\cftdotfill{\cftdotsep}}
\renewcommand{\cftsubsubsecleader}{\cftdotfill{\cftdotsep}}
\renewcommand{\cftsecfont}{\fontsize{12}{14}\bfseries}
\renewcommand{\cftsecpagefont}{\fontsize{12}{14}\bfseries}
\renewcommand{\cftsecaftersnum}{.\enspace}
\settowidth{\cftsecnumwidth}{VIII.\enspace}
\renewcommand{\cftsubsecaftersnum}{.}
\renewcommand{\cftsubsubsecaftersnum}{.}

% --- Надписи: таблици ОТГОРЕ, фигури ОТДОЛУ (А.8, А.9) ------------------
\usepackage{caption}
\usepackage{subcaption}
\DeclareCaptionLabelFormat{tablelabel}{Таблица #2}
\captionsetup[table]{labelformat=tablelabel, labelsep=period, position=above,
                     font={bf,small}, justification=centering, singlelinecheck=true}
\DeclareCaptionLabelFormat{figurelabel}{Фигура #2}
\captionsetup[figure]{labelformat=figurelabel, labelsep=period, position=below,
                      font={bf,small}, justification=centering, singlelinecheck=true}

% Препратки в текста: „(Табл. 3)“ / „(Фиг. 7)“ — задължителни по А.8/А.9.
\newcommand{\tabref}[1]{(Табл.~\ref{#1})}
\newcommand{\figref}[1]{(Фиг.~\ref{#1})}

% --- Математика ---------------------------------------------------------
\usepackage{amsmath,amssymb,amsfonts}
\usepackage{amsthm}

% --- Графики ------------------------------------------------------------
\usepackage{graphicx}
\usepackage{float}
\usepackage{tikz}
\usetikzlibrary{decorations.pathreplacing,arrows.meta,positioning,shapes.geometric,fit,calc}
\usepackage{adjustbox}
\usepackage{pdfpages}     % за подвързаните задание и декларация, вж. Б.5

% --- Таблици ------------------------------------------------------------
\usepackage{array}
% Ляво подравнена колона с фиксирана ширина. Тесните `p{}` колони с ДВУСТРАННО
% подравняване (наследено от основния текст) разкъсват редовете с огромни
% междудумни интервали — правилото за justified от А.4 се отнася за ОСНОВНИЯ текст,
% не за клетки от 3 cm. Всички таблици с текстови клетки ползват L.
\newcolumntype{L}[1]{>{\RaggedRight\arraybackslash}p{#1}}
\usepackage{booktabs}
\usepackage[table]{xcolor}% \cellcolor за топлинната карта
\usepackage{longtable}    % „продължение“ при пренасяне на таблица (А.8)
\usepackage{multirow}
\usepackage{tabularx}

% --- Помощни ------------------------------------------------------------
\usepackage{hyperref}
\usepackage{xurl}
\usepackage{enumitem}
\usepackage{textcomp}
\usepackage{listings}
\urlstyle{same}
% КРИТИЧНО: `breaklines` на listings вмъква ДИСКРЕЦИОННИ прекъсвания, а глобалните
% \hyphenpenalty/\exhyphenpenalty=10000 (правилото „без пренасяния“, А.4) ги правят
% безкрайно скъпи — тогава дългите редове в листингите НЕ се пренасят и излизат извън
% полето. Проверено изолирано: с наказанията 1 препълване, без тях 0. Наказанията се
% нулират САМО вътре в листинг, така че текстът си остава без пренасяния.
\makeatletter
\lst@AddToHook{PreInit}{\hyphenpenalty=50\exhyphenpenalty=50}
\makeatother
\renewcommand{\lstlistingname}{Листинг}
\renewcommand{\lstlistlistingname}{Списък на листингите}

% Italic САМО за ключови понятия (А.4) — семантичен макрос, за да е ревизуемо.
\newcommand{\term}[1]{\emph{#1}}
% Инлайн код/идентификатори.
\DeclareRobustCommand{\code}[1]{\texttt{#1}}
% Глифи за таблиците с оценки.
\newcommand{\cmark}{\ensuremath{\checkmark}}
\newcommand{\xmark}{\ensuremath{\times}}

% Заснемане на екранна форма с тънка черна рамка. Рамката НЕ е украса: екранните
% форми на приложението са върху почти бял фон и се сливат с бялата страница, така
% че без нея не се вижда къде свършва изображението и започва листът — а размерът
% на изображението е част от твърдението. \fboxsep=0pt прилепва рамката към
% пиксела; 0.4 pt е най-тънката, която се възпроизвежда надеждно при печат.
\newcommand{\capture}[2][0.44\textwidth]{%
  {\setlength{\fboxsep}{0pt}\setlength{\fboxrule}{0.4pt}%
   \fbox{\includegraphics[width=#1]{#2}}}}

% Настолна екранна форма през цялата ширина на текста. Заснемането е 2880x1800 px,
% тоест при 15,7 cm ширина плътността е около 460 dpi и надписите в интерфейса се
% четат. По-малък размер ги превръща в сиви петна — това е причината екранните форми
% да НЕ се редят по две на ред.
%   #1 ключ на формата  #2 описание
% NB: рамката добавя 2 x 0,4 pt към ширината, затова изображението се задава с
% \textwidth минус дебелината на рамката — иначе редът е препълнен с точно 0,8 pt.
\newcommand{\screenshot}[2]{%
  \begin{figure}[htbp]\centering
    \capture[\dimexpr\textwidth-0.8pt\relax]{images/screens/#1_1440.png}
    \caption{#2}\label{fig:ss-#1}
  \end{figure}}

% Трите форми при ширина 375 px една до друга: доказателството за поведението при
% тесен екран, събрано на едно място вместо разпръснато из целия раздел.
\newcommand{\phonestrip}[1]{%
  \begin{figure}[htbp]\centering
    \begin{tabular}{@{}ccc@{}}
      \small\textbf{Калкулатор} & \small\textbf{Проекти} & \small\textbf{Документация} \\
      \capture[0.30\textwidth]{images/screens/calculator_375.png} &
      \capture[0.30\textwidth]{images/screens/dashboard_375.png} &
      \capture[0.30\textwidth]{images/screens/docs_375.png} \\
    \end{tabular}
    \caption{#1}\label{fig:phonestrip}
  \end{figure}}

% Съпоставка на една и съща форма при две ширини. Двете заснемания имат много
% различно съотношение (16:10 срещу 1:2,17), затова се подравняват по ОБЩА ВИСОЧИНА,
% а не по зададени поотделно ширини: иначе телефонът изяжда височината на реда, а
% настолната форма се смалява до нечетливост. При обща височина H настолната форма е
% 1,6H широка, а телефонната 0,46H, тоест 2,06H <= \textwidth дава H = 0,485\textwidth.
%   #1 ключ на формата  #2 описание
\newcommand{\viewportpair}[2]{%
  \begin{figure}[htbp]\centering
    \begin{tabular}{@{}cc@{}}
      \small\textbf{1440 px (настолен)} & \small\textbf{375 px (телефон)} \\
      \capture[0.755\textwidth]{images/screens/#1_1440.png} &
      \capture[0.218\textwidth]{images/screens/#1_375.png} \\
    \end{tabular}
    \caption{#2}\label{fig:vp-#1}
  \end{figure}}

% Видим маркер за незапълнено място. Търси се с: grep -rn 'TODO' chapters/
\newcommand{\TODO}[1]{\textcolor{red}{\textbf{[TODO: #1]}}}

\tikzset{
    block/.style={draw, rounded corners, fill=blue!5, align=center, font=\small,
        minimum width=3.2cm, minimum height=1.05cm, inner sep=6pt},
    accent/.style={draw, rounded corners, fill=green!8, align=center, font=\small,
        minimum width=3.0cm, minimum height=0.95cm, inner sep=6pt},
    danger/.style={draw, rounded corners, fill=red!6, align=center, font=\small,
        minimum width=3.0cm, minimum height=0.95cm, inner sep=6pt},
    note/.style={draw, dashed, rounded corners, fill=gray!8, align=left, font=\small,
        inner sep=6pt},
    line/.style={-Latex, thick}
}

% --- Единично разстояние в нетекстови среди -----------------------------
\AtBeginEnvironment{tikzpicture}{\setstretch{1.1}}
\AtBeginEnvironment{lstlisting}{\setstretch{1.1}}
\AtBeginEnvironment{verbatim}{\setstretch{1.1}}
\AtBeginEnvironment{tabular}{\setstretch{1.1}}
\AtBeginEnvironment{longtable}{\setstretch{1.1}}

\BeforeBeginEnvironment{tikzpicture}{\begin{adjustbox}{max width=\textwidth}}
\AfterEndEnvironment{tikzpicture}{\end{adjustbox}}

% --- Оцветяване на изходния текст ---------------------------------------
% Кодовата база е на JavaScript (Node.js, браузър, React Native), затова
% основният език за листингите е JavaScript. Диалектите (JSON, CSS, HTML) се
% задават локално с \begin{lstlisting}[language=…].
\definecolor{cppblue}{RGB}{0,0,180}
\definecolor{cppgreen}{RGB}{0,128,0}
\definecolor{cppred}{RGB}{180,0,0}
\definecolor{cppgray}{RGB}{128,128,128}

% `listings` не носи описание на JavaScript, затова езикът се дефинира тук.
\lstdefinelanguage{JavaScript}{
    keywords={await, async, break, case, catch, class, const, continue, debugger,
              default, delete, do, else, export, extends, finally, for, function,
              if, import, in, instanceof, let, new, of, return, super, switch,
              this, throw, try, typeof, var, void, while, with, yield,
              true, false, null, undefined},
    sensitive=true,
    morecomment=[l]{//},
    morecomment=[s]{/*}{*/},
    morestring=[b]',
    morestring=[b]",
    morestring=[b]`,
}

% `listings` не носи описание и на CSS.
\lstdefinelanguage{CSS}{
    morekeywords={display, flex, grid, width, min-width, max-width, margin, padding,
                  position, overflow, overflow-x, align-items, justify-content,
                  grid-template-columns, grid-column, grid-row, top, gap},
    sensitive=true,
    morecomment=[s]{/*}{*/},
    morestring=[b]',
    morestring=[b]",
}

\lstdefinelanguage{JSONlang}{
    keywords={true,false,null},
    sensitive=true,
    morestring=[b]",
    morecomment=[l]{//},
}

\lstset{
    language=JavaScript,
    basicstyle=\footnotesize\ttfamily,
    % Пренасяне на дълги редове. `breakatwhitespace=false` е СЪЩЕСТВЕНО: коментарите в
    % изходния текст стигат 110 знака без удобно място за прекъсване, и при `true` редът
    % изобщо не се пренася, а излиза извън полето. `linewidth` се задава изрично, защото
    % рамката и левият отстъп иначе не се приспадат от наличната ширина.
    breaklines=true, breakatwhitespace=false, breakautoindent=false, breakindent=0pt,
    linewidth=\dimexpr\linewidth-2.5em\relax,
    columns=flexible, keepspaces=true,
    postbreak=\mbox{\textcolor{cppgray}{\ensuremath{\hookrightarrow}}\space},
    xleftmargin=1.5em, xrightmargin=0.5em,
    keywordstyle=\color{cppblue}\bfseries,
    stringstyle=\color{cppred},
    commentstyle=\color{cppgreen},
    identifierstyle=\color{black},
    showstringspaces=false,
    numbers=left, numberstyle=\tiny\color{cppgray}, numbersep=8pt,
    aboveskip=1.5em, belowskip=1.5em,
    frame=single, framerule=0.5pt, rulecolor=\color{cppgray},
    tabsize=4, captionpos=b,
    morekeywords={const, let, async, await, of, class, extends, export, import,
                  from, static, yield, require, module, typeof, instanceof},
    emph={app, router, express, res, req, next, fetch, localStorage, document,
          window, WTApi, WTAuth, WALLTOPIA\_LOADS, renderResults, clampAndRender},
    emphstyle=\color{cppblue},
    inputencoding=utf8,
    % Не-ASCII знаци, срещащи се в коментарите на изходния текст. Без тези редове
    % \lstinputlisting отказва с „Invalid UTF-8 byte sequence“ — listings не разбира
    % многобайтов вход, който не е в literate таблицата.
    literate={—}{{---}}1 {→}{{$\rightarrow$}}1 {×}{{$\times$}}1 {…}{{...}}1 {§}{{\S}}1
{а}{{\cyra}}1 {б}{{\cyrb}}1 {в}{{\cyrv}}1 {г}{{\cyrg}}1 {д}{{\cyrd}}1
             {е}{{\cyre}}1 {ж}{{\cyrzh}}1 {з}{{\cyrz}}1 {и}{{\cyri}}1 {й}{{\cyrishrt}}1
             {к}{{\cyrk}}1 {л}{{\cyrl}}1 {м}{{\cyrm}}1 {н}{{\cyrn}}1 {о}{{\cyro}}1
             {п}{{\cyrp}}1 {р}{{\cyrr}}1 {с}{{\cyrs}}1 {т}{{\cyrt}}1 {у}{{\cyru}}1
             {ф}{{\cyrf}}1 {х}{{\cyrh}}1 {ц}{{\cyrc}}1 {ч}{{\cyrch}}1 {ш}{{\cyrsh}}1
             {щ}{{\cyrshch}}1 {ъ}{{\cyrhrdsn}}1 {ь}{{\cyrsftsn}}1 {ю}{{\cyryu}}1 {я}{{\cyrya}}1
             {А}{{\CYRA}}1 {Б}{{\CYRB}}1 {В}{{\CYRV}}1 {Г}{{\CYRG}}1 {Д}{{\CYRD}}1
             {Е}{{\CYRE}}1 {Ж}{{\CYRZH}}1 {З}{{\CYRZ}}1 {И}{{\CYRI}}1 {Й}{{\CYRISHRT}}1
             {К}{{\CYRK}}1 {Л}{{\CYRL}}1 {М}{{\CYRM}}1 {Н}{{\CYRN}}1 {О}{{\CYRO}}1
             {П}{{\CYRP}}1 {Р}{{\CYRR}}1 {С}{{\CYRS}}1 {Т}{{\CYRT}}1 {У}{{\CYRU}}1
             {Ф}{{\CYRF}}1 {Х}{{\CYRH}}1 {Ц}{{\CYRC}}1 {Ч}{{\CYRCH}}1 {Ш}{{\CYRSH}}1
             {Щ}{{\CYRSHCH}}1 {Ъ}{{\CYRHRDSN}}1 {Ь}{{\CYRSFTSN}}1 {Ю}{{\CYRYU}}1 {Я}{{\CYRYA}}1,
}
